\section{Future Work}
\label{sec:future}

\subsection{Instruction-set coverage}

The most direct limitation is the x86-64 subset. Future work includes:

\begin{itemize}
  \item general memory addressing and writable data segments with a private
        copy and explicit synchronization policy;
  \item indirect jumps and jump tables with a validated dispatch table;
  \item AVX2/AVX-256 and selected floating-point operations;
  \item a broader SSE/AVX semantic set, including shuffles and conversions;
  \item atomic and fence semantics with a documented memory model.
\end{itemize}

Each extension must preserve the fail-closed contract: an unsupported form is a
build error, and a supported form must have an exact semantic model.

\subsection{External calls and language exceptions}

A selective virtualizer that can call external functions needs a host-call stub
mechanism with an explicit ABI boundary, a return-value marshaling policy, and a
stack-unwind model. Language exceptions require a native unwind frame or a
guest-level exception representation. These are substantial design changes, not
minor decoder additions.

\subsection{Platform coverage}

The Linux path is the primary release gate. Future platform work includes a
fuller Win64 runtime, Mach-O support, ARM64, and a 32-bit x86 backend. Each
platform has different unwind metadata, section geometry, ABI details, and
anti-debug surfaces; a platform port is not a simple recompile.

\subsection{Anti-debug and integrity}

The hardened tier currently provides user-mode tracer detection, PEB checks,
RDTSC/LFENCE, debug-register reporting, and linked text checksums. Future work
includes:

\begin{itemize}
  \item a formally reviewed hardened boundary with a reduced unsafe surface;
  \item kernel-assisted integrity and anti-debug for deployments that can
        require it;
  \item signed artifacts and a verified update chain;
  \item a memory-encryption or page-protection design that does not rely on
        unsafe access in the safe tier.
\end{itemize}

\subsection{Constant-time and side channels}

The fixed-pass mode can be extended with a more precise convergence analysis, a
smaller memory-scan cost, and a formal model of the source-level property. A
separate line of work is required for speculative execution, cache-bank
conflicts, and power/EM leakage. Those are not addressed by branchless bytecode
alone.

\subsection{MBA and diversity}

The MBA layer can be improved with:

\begin{itemize}
  \item a larger library of exact non-linear identities;
  \item equivalence checking over a bounded bit-vector model at build time;
  \item per-block MBA shape selection rather than one wrapper per function;
  \item a differential study against newer simplifiers and SMT-based
        deobfuscators.
\end{itemize}

Diversity can be improved by mutating the interpreter implementation itself, not
only its dispatch tables and layout. That requires a careful balance between code
size, compiler canonicalization, and verification cost.

\subsection{Formal methods}

The current proofs cover bounded helpers and boundary models. A future direction
is a mechanized model of the virtual ISA and a refinement proof from the lifter
to the interpreter for the supported subset. The decoder is the largest obstacle:
a verified decoder for even a small x86 subset is a substantial project. A
smaller-step approach is to verify the encoder/decoder round-trip for the virtual
ISA and to keep the native decoder under differential and fuzz testing.

\subsection{Independent evaluation}

The most important non-technical gap is independent review. The repository
contains reproducible evidence, but the authors are also the implementers. An
external audit should examine the custom payload construction, the unsafe
hardened boundary, the decoder's fail-closed behavior, and the Windows unwind
metadata. A public challenge with a fixed artifact and a time budget would be a
useful complement to the internal evidence.
